\chapter{基础算法与数据结构}

\section{前缀和与差分}

\cardmeta{基础必会}{区间静态求和、区间统一加}{预处理 \(O(n)\)，查询/修改 \(O(1)\)}{1-based 或 0-based 必须在卡片中统一}

\begin{lstlisting}[style=cpp]
// pre[i] = a[1..i] 的和；下标 0 留作空前缀
vector<int> pre(n + 1, 0);
for (int i = 1; i <= n; ++i) pre[i] = pre[i - 1] + a[i];
auto range_sum = [&](int l, int r) {
    // 闭区间 [l,r] 的和 = 两个前缀相减
    return pre[r] - pre[l - 1];
};

// diff[l..r] 全部加 v，只需要修改两个端点
vector<int> diff(n + 2, 0);
auto add = [&](int l, int r, int v) {
    diff[l] += v;
    diff[r + 1] -= v;
};
for (int i = 1; i <= n; ++i) {
    // 对差分做前缀和，还原每个位置最终增加了多少
    diff[i] += diff[i - 1];
    a[i] += diff[i];
}
\end{lstlisting}

\begin{card}
\textbf{完整样例}\\
输入：
\begin{lstlisting}
4 2
1 2 3 4
1 3 5
2 4 -1
\end{lstlisting}
输出：
\begin{lstlisting}
6 12 12 3
\end{lstlisting}
\end{card}

\section{并查集}

\cardmeta{基础必会}{连通块、合并集合、最小生成树}{单次 \(O(\alpha(n))\)}{节点编号 1 到 \(n\)}

\begin{lstlisting}[style=cpp]
struct DSU {
    // fa 是父节点；sz 只在根节点上有意义，用于按大小合并
    vector<int> fa, sz;
    explicit DSU(int n = 0) { init(n); }
    void init(int n) {
        fa.resize(n + 1);
        sz.assign(n + 1, 1);
        iota(fa.begin(), fa.end(), 0);
    }
    int find(int x) {
        // 路径压缩：查询后把 x 直接挂到根，降低后续复杂度
        return fa[x] == x ? x : fa[x] = find(fa[x]);
    }
    bool unite(int x, int y) {
        x = find(x); y = find(y);
        // 同根说明已经连通，不能重复合并
        if (x == y) return false;
        // 小树挂到大树，避免并查树退化
        if (sz[x] < sz[y]) swap(x, y);
        fa[y] = x; sz[x] += sz[y];
        return true;
    }
};
\end{lstlisting}

\begin{card}
\textbf{样例：}\\
输入：
\begin{lstlisting}
5 4
1 2
2 3
4 5
1 3
\end{lstlisting}
输出：
\begin{lstlisting}
3
\end{lstlisting}
表示最后有 3 个连通块。
\end{card}

\section{树状数组}

\cardmeta{现场常用}{单点加、前缀和、逆序对、离散化后计数}{单次 \(O(\log n)\)}{下标从 1 开始}

\begin{lstlisting}[style=cpp]
struct Fenwick {
    // bit[x] 保存以 x 结尾、长度 lowbit(x) 的区间和
    int n;
    vector<int> bit;
    explicit Fenwick(int n = 0) { init(n); }
    void init(int n_) { n = n_; bit.assign(n + 1, 0); }
    void add(int x, int v) {
        // x += lowbit(x)：更新所有覆盖位置 x 的节点
        for (; x <= n; x += x & -x) bit[x] += v;
    }
    int sum(int x) const {
        int res = 0;
        // x -= lowbit(x)：拆成互不重叠的若干块
        for (; x > 0; x -= x & -x) res += bit[x];
        return res;
    }
    int range_sum(int l, int r) const {
        return sum(r) - sum(l - 1);
    }
};
\end{lstlisting}

\begin{warnbox}
树状数组不能直接维护任意区间最值。若要区间加、单点查，把加法作用在差分数组上；若要区间加、区间和，需要两棵树或改用线段树。
\end{warnbox}

\section{区间加、区间和线段树}

\cardmeta{现场常用}{区间加、区间和、懒标记}{建树 \(O(n)\)，操作 \(O(\log n)\)}{节点池按 \(4n\) 开}

\begin{lstlisting}[style=cpp]
struct SegTree {
    // sum[p] 是节点 p 对应区间的总和，lazy[p] 是尚未下传的区间加
    int n;
    vector<int> sum, lazy;
    explicit SegTree(int n = 0) { init(n); }
    void init(int n_) {
        n = n_;
        sum.assign(4 * n + 5, 0);
        lazy.assign(4 * n + 5, 0);
    }
    void apply(int p, int l, int r, int v) {
        // 整段加 v：总和增加 v * 区间长度，标记累加
        sum[p] += v * (r - l + 1);
        lazy[p] += v;
    }
    void push(int p, int l, int r) {
        // 叶子没有子节点；没有标记也无需下传
        if (!lazy[p] || l == r) return;
        int mid = (l + r) >> 1;
        apply(p << 1, l, mid, lazy[p]);
        apply(p << 1 | 1, mid + 1, r, lazy[p]);
        lazy[p] = 0;
    }
    void add(int p, int l, int r, int ql, int qr, int v) {
        // 当前区间被完全覆盖时，打标记后立即返回
        if (ql <= l && r <= qr) return apply(p, l, r, v);
        // 部分覆盖前必须先把父标记推给孩子
        push(p, l, r);
        int mid = (l + r) >> 1;
        if (ql <= mid) add(p << 1, l, mid, ql, qr, v);
        if (qr > mid) add(p << 1 | 1, mid + 1, r, ql, qr, v);
        sum[p] = sum[p << 1] + sum[p << 1 | 1];
    }
    int query(int p, int l, int r, int ql, int qr) {
        if (ql <= l && r <= qr) return sum[p];
        push(p, l, r);
        int mid = (l + r) >> 1, res = 0;
        if (ql <= mid) res += query(p << 1, l, mid, ql, qr);
        if (qr > mid) res += query(p << 1 | 1, mid + 1, r, ql, qr);
        return res;
    }
};
\end{lstlisting}

\begin{tipbox}
从 `muban.cpp` 复制后，先把初始数组逐点调用 `add(1,1,n,i,i,a[i])`，再接题目操作。每次访问子区间前必须 `push`。
\end{tipbox}

\section{二分答案}

\cardmeta{基础必会}{最小的最大值、最大的最小值、可行性判定}{\(O(n\log V)\)}{先确认判定函数单调}

\begin{lstlisting}[style=cpp]
int lo = *max_element(a.begin(), a.end()); // 答案下界：至少容纳单个最大值
int hi = accumulate(a.begin(), a.end(), 0LL); // 答案上界：全部元素放一组
while (lo < hi) {
    int mid = (lo + hi) >> 1;
    // check(mid) 必须是单调函数：可行则尝试更小，否则只能增大
    if (check(mid)) hi = mid;
    else lo = mid + 1;
}
cout << lo << endl;
\end{lstlisting}

\begin{warnbox}
二分答案不是看到“最小/最大”就能套。必须先证明 `check(x)` 的真假随 \(x\) 单调变化，且 `lo`、`hi` 覆盖真实答案。
\end{warnbox}
